\section{Kontextfreie Grammatiken}

\subsection{Definition}
Eine kontextfreie Grammatik (KFG) ist ein Tupel
\[
	G = (N, \Sigma, P, S)
\]
mit
\begin{itemize}
	\item Alphabet der \textbf{Nichtterminalen} $N$,
	\item Alphabet der \textbf{Terminalen} $\Sigma$ mit $N \cap \Sigma = \emptyset$,
	\item Endliche Menge \textbf{Produktionen} $P \subseteq N \times (N \cup \Sigma)^*$,
	\item \textbf{Startsymbol} $S \in N$.
\end{itemize}

Eine Produktion hat die Form
\[
	A \to \alpha \quad \text{mit } A \in N,\; \alpha \in (N \cup \Sigma)^*.
\]
\(\alpha \in (N \cup \Sigma)^*\) wird auch als \textbf{Satzformel} bezeichnet.

\subsection{Ableitung}
Der Ableitungsschritt ist definiert durch
\[
	uAv \Rightarrow_G u\alpha v
\]
genau dann, wenn $A \to \alpha \in P$ und $u,v \in (N \cup \Sigma)^*$.

Die reflexiv-transitive Hülle wird mit $\Rightarrow_G^*$ bezeichnet.

\subsection{Erzeugte Sprache}
Die von $G$ erzeugte Sprache ist
\[
	L(G) = \{\,w \in \Sigma^* \mid S \Rightarrow_G^* w\,\}.
\]
Eine Sprache $L \subseteq \Sigma^*$ heisst kontextfrei, falls es eine KFG $G$ mit $L=L(G)$ gibt.

\subsection{Links- und Rechtsableitung}
\begin{itemize}
	\item \textbf{Linksableitung:} In jedem Schritt wird das linkeste Nichtterminal ersetzt.
	\item \textbf{Rechtsableitung:} In jedem Schritt wird das rechteste Nichtterminal ersetzt.
\end{itemize}
Beide beschreiben dieselbe erzeugte Sprache.

\subsection{Ableitungsbaum}
Ein Ableitungsbaum repräsentiert die hierarchische Struktur einer Ableitung:
\begin{itemize}
	\item Wurzel ist $S$.
	\item Innere Knoten sind Nichtterminale.
	\item Blätter sind Terminalsymbole.
\end{itemize}
Die von links nach rechts gelesenen Blätter ergeben das erzeugte Wort.

\subsection{Mehrdeutigkeit}
Eine Grammatik ist mehrdeutig, wenn es ein Wort $w \in L(G)$ mit mindestens zwei verschiedenen Ableitungsbäumen gibt.
Eine Sprache heisst inherent mehrdeutig, wenn jede Grammatik, die sie erzeugt, mehrdeutig ist.
